% Dissolution du chlorure de sodium dans l'eau : dissociation et hydratation des ions (L'essentiel, chapitre 7)
\begin{tikzpicture}[x=1cm,y=1cm,
  aO/.style={circle,fill=otRed!85,minimum size=3.8mm,inner sep=0pt},
  aH/.style={circle,fill=white,draw=otInk!55,line width=0.4pt,minimum size=2.4mm,inner sep=0pt}]
\newcommand{\eau}[3]{%
  \draw[otInk!70,line width=0.7pt] ($(#1,#2)+(217.75+#3:0.3)$) -- (#1,#2) -- ($(#1,#2)+(322.25+#3:0.3)$);
  \node[aH] at ($(#1,#2)+(217.75+#3:0.3)$) {}; \node[aH] at ($(#1,#2)+(322.25+#3:0.3)$) {};
  \node[aO] at (#1,#2) {};}
% Cristal
\foreach \i in {0,...,3}{\foreach \j in {0,...,3}{
  \pgfmathtruncatemacro\parite{mod(\i+\j,2)}
  \ifnum\parite=0 \shade[ball color=otGreen!80] ({0.6*\i},{0.6*\j}) circle (0.2);
  \else \shade[ball color=otNavy!70] ({0.6*\i},{0.6*\j}) circle (0.11); \fi}}
\node[pcLegendeG,anchor=north] at (0.9,-0.4) {cristal ionique};
\draw[otGreenDark,line width=1pt,-{Stealth[length=2.5mm]}] (2.4,0.9) -- (3.6,0.9) node[midway,above,pcLegende,text=otGreenDark] {eau};
% Ion sodium hydraté : les atomes d'oxygène pointent vers l'ion
\begin{scope}[shift={(5.0,0.9)}]
  \foreach \p in {0,72,...,288}{\eau{{0.62*cos(\p)}}{{0.62*sin(\p)}}{\p-270}}
  \shade[ball color=otNavy!70] (0,0) circle (0.15); \node[font=\otSemi\fontsize{6}{7}\selectfont,text=white] at (0,0) {+};
  \node[pcLegendeG,anchor=north] at (0,-1.3) {ion sodium hydraté};
\end{scope}
% Ion chlorure hydraté : un atome d'hydrogène pointe vers l'ion
\begin{scope}[shift={(8.2,0.9)}]
  \foreach \p in {0,72,...,288}{\eau{{0.88*cos(\p)}}{{0.88*sin(\p)}}{\p-37.75}}
  \shade[ball color=otGreen!80] (0,0) circle (0.27); \node[font=\otSemi\fontsize{7}{8}\selectfont,text=white] at (0,0) {$-$};
  \node[pcLegendeG,anchor=north] at (0,-1.3) {ion chlorure hydraté};
\end{scope}
\end{tikzpicture}
\par\smallskip\legende{L'eau dissocie le cristal, puis ses molécules polaires entourent chaque ion : pôle $\delta^{-}$ (oxygène) vers les cations, pôle $\delta^{+}$ (hydrogène) vers les anions.}
